\documentclass[a4paper, 11pt]{article}

\usepackage[T1]{fontenc}
\usepackage{fix-cm}
\usepackage{amsmath}
\usepackage{courier}
\usepackage[scaled=0.85]{beramono}
\usepackage{lmodern}

\usepackage[polish]{babel}

\usepackage{etoolbox}
\preto{\enumerate}{\par\nobreak}

\usepackage[
  top=2.5cm, bottom=2.5cm,
  left=2.8cm, right=2.8cm,
  headheight=22pt
]{geometry}

\usepackage{tocloft}
\setlength{\cftbeforesecskip}{0.3em}
\setlength{\cftbeforesubsecskip}{0pt}

\usepackage{xcolor}
\definecolor{primary}{HTML}{1A1A1A}
\definecolor{accent}{HTML}{1A1A1A}
\definecolor{lightbg}{HTML}{F0F0F0}
\definecolor{codebg}{HTML}{EBEBEB}
\definecolor{codefg}{HTML}{1A1A1A}
\definecolor{rulegray}{HTML}{999999}
\definecolor{partcolor}{HTML}{1E3A5F}

% Kolory dla popup'ów (z dokumentacji funkcjonalnej)
\definecolor{errorbg}{HTML}{FDECEA}
\definecolor{errorframe}{HTML}{C0392B}
\definecolor{warnbg}{HTML}{FEF9E7}
\definecolor{warnframe}{HTML}{F39C12}

\usepackage[nobottomtitles*]{titlesec}
\titleformat{\section}
  {\color{primary}\large\bfseries}
  {\color{accent}\thesection.}{0.6em}{}[{\color{rulegray}\titlerule[0.6pt]}]
\titleformat{\subsection}
  {\color{primary}\normalsize\bfseries}
  {\color{accent}\thesubsection.}{0.5em}{}
\titlespacing{\section}{0pt}{1.4em}{0.7em}
\titlespacing{\subsection}{0pt}{1.0em}{0.4em}

\usepackage{fancyhdr}
\pagestyle{fancy}
\fancyhf{}
\fancyhead[L]{\small\color{primary}\textbf{Wizualizacja grafu planarnego}}
\fancyhead[R]{\small\color{rulegray}Dokumentacja końcowa -- Java}
\fancyfoot[C]{\small\color{rulegray}\thepage}

\usepackage{enumitem}
\setlist[itemize]{
  leftmargin=1.8em,
  itemsep=0.3em,
  topsep=0.4em
}
\setlist[enumerate]{
  leftmargin=1.8em,
  itemsep=0.3em,
  topsep=0.4em
}

\usepackage{tcolorbox}
\tcbuselibrary{skins, breakable}

\newtcolorbox{cmdbox}{
  enhanced,
  colback=codebg,
  colframe=primary,
  arc=3pt,
  boxrule=0.8pt,
  left=8pt, right=8pt, top=5pt, bottom=5pt,
  fontupper=\ttfamily\small\color{codefg}
}

\newtcolorbox{filebox}{
  enhanced,
  colback=lightbg,
  colframe=rulegray,
  arc=3pt,
  boxrule=0.6pt,
  left=10pt, right=10pt, top=6pt, bottom=6pt,
  fontupper=\ttfamily\small\color{primary}
}

\newtcolorbox{errorbox}[1][]{
  enhanced,
  colback=errorbg,
  colframe=errorframe,
  arc=3pt,
  boxrule=1pt,
  left=8pt, right=8pt, top=5pt, bottom=5pt,
  title={\bfseries\color{errorframe}\textbullet\ B\l{}\'ad krytyczny #1},
  fontupper=\ttfamily\small\color{codefg}
}

\newtcolorbox{warnbox}[1][]{
  enhanced,
  colback=warnbg,
  colframe=warnframe,
  arc=3pt,
  boxrule=1pt,
  left=8pt, right=8pt, top=5pt, bottom=5pt,
  title={\bfseries\color{warnframe}\textbullet\ Ostrze\.zenie #1},
  fontupper=\ttfamily\small\color{codefg}
}

\usepackage[hidelinks]{hyperref}
\usepackage{needspace}
\usepackage{booktabs}
\usepackage{array}
\usepackage{titletoc}
\usepackage{graphicx}
\usepackage{float}

% Pakiety dla schematu GUI
\usepackage{tikz}
\usetikzlibrary{shapes.geometric, arrows.meta, positioning, fit, backgrounds}

\setlength{\parskip}{0.4em}
\setlength{\parindent}{0pt}
\setlength{\emergencystretch}{3em}

% ---- Styl strony rozdziałowej dla części ----
\newcommand{\partpage}[2]{%
  \clearpage
  \thispagestyle{empty}
  \vspace*{6cm}
  \begin{center}
    {\color{rulegray}\rule{0.4\textwidth}{0.6pt}}\\[1em]
    {\fontsize{11}{13}\selectfont\scshape\color{rulegray}#1}\\[0.6em]
    {\fontsize{28}{34}\selectfont\bfseries\color{primary}#2}\\[1em]
    {\color{rulegray}\rule{0.4\textwidth}{0.6pt}}
  \end{center}
  \clearpage
}

\begin{document}

% ============================================================
% STRONA TYTUŁOWA
% ============================================================
\begin{titlepage}
  \centering
  \vspace*{4cm}

  {\fontsize{32}{38}\selectfont Wizualizacja grafu\\planarnego}

  \vspace{0.8cm}

  {\large\scshape\color{rulegray} Dokumentacja końcowa --- część Java}

  \vspace{1.2cm}

  {\small\scshape Wiktor Godula\\ Przemysław Feliniak}

  \vfill

  {\small\color{rulegray} \today}
\end{titlepage}

% ============================================================
% SPIS TREŚCI
% ============================================================
\tableofcontents
\newpage

% ============================================================
% ============================================================
%   CZĘŚĆ I: DOKUMENTACJA FUNKCJONALNA
% ============================================================
% ============================================================

\partpage{Część I}{Dokumentacja funkcjonalna}

% =======================================================
\section{Wst\k{e}p}
% =======================================================

Celem projektu jest stworzenie interaktywnej aplikacji okienkowej w j\k{e}zyku Java (technologia Swing), stanowi\k{a}cej uzupe\l{}nienie cz\k{e}\'sci obliczeniowej zrealizowanej w j\k{e}zyku C. Aplikacja wczytuje wyniki dzia\l{}ania programu C (graf planarny wraz z wyliczonymi wsp\'o\l{}rz\k{e}dnymi) i pozwala na ich komfortow\k{a} wizualizacj\k{e}.

Rola wzgl\k{e}dem komponentu~C.
Aplikacja Java pe\l{}ni rol\k{e} graficznej nakładki steruj\k{a}cej i wizualizuj\k{a}cej wzgl\k{e}dem samodzielnego programu konsolowego \texttt{graf\_planarny} napisanego w j\k{e}zyku~C. Silnik obliczeniowy (plik wykonywalny \texttt{graf\_planarny}) musi by\'c skompilowany i dost\k{e}pny w podkatalogu \texttt{core/} katalogu roboczego aplikacji (np. \texttt{core/graf\_planarny} na Linux/macOS lub \texttt{core/graf\_planarny.exe} na Windows). Aplikacja nie przeszukuje innych ścieżek.

Aplikacja Java zawsze uruchamia \texttt{graf\_planarny} jako osobny proces. Na podstawie wyb\'or\'ow u\.zytkownika dokonanych w Pasku Narz\k{e}dziowym (plik wej\'sciowy z list\k{a} kraw\k{e}dzi, algorytm, liczba iteracji, format zapisu) aplikacja konstruuje i uruchamia polecenie zgodne ze specyfikacj\k{a} programu C opisan\k{a} w dokumentacji ko\'ncowej j\k{e}zyka~C (sekcja~3.1), np.:

\begin{cmdbox}
./graf\_planarny -i wejscie.txt -o wynik\_tmp.txt -a fr -n 600
\end{cmdbox}

Wynik zapisywany jest do pliku tymczasowego, kt\'ory aplikacja nast\k{e}pnie parsuje i wy\'swietla w obszarze roboczym. Aplikacja Java nie zawiera w\l{}asnej implementacji algorytm\'ow obliczeniowych --- ca\l{}a logika matematyczna pozostaje wy\l{}\k{a}cznie po stronie programu~C. Aby wywo\l{}anie si\k{e} powiod\l{}o, \'scie\.zka do pliku wykonywalnego \texttt{graf\_planarny} musi by\'c dost\k{e}pna (domy\'slnie: katalog roboczy aplikacji lub zmienna \'srodowiskowa \texttt{PATH}).

Program dostarcza narz\k{e}dzi do wy\'swietlania wczytanych graf\'ow, modyfikacji sposobu ich prezentacji, a tak\.ze manualnej zmiany pozycji w\k{e}z\l{}\'ow i uruchamiania algorytm\'ow przez program C. Dokument ten jest jednocze\'snie instrukcj\k{a} dla u\.zytkownika, specyfikacj\k{a} funkcjonaln\k{a} i opisem architektury przeznaczonym dla zespo\l{}u programistycznego.

% =======================================================
\section{Zestawienie g\l{}ównych funkcjonalno\'sci}
% =======================================================

Poni\.zsze punkty przedstawiaj\k{a} kluczowe mo\.zliwo\'sci aplikacji z punktu widzenia u\.zytkownika:

\begin{itemize}
  \item \textbf{Wczytywanie grafu} --- dwa tryby: \textit{Tekstowy} wczytuje plik z listą krawędzi, \textit{Binarny} wczytuje plik ze współrzędnymi wierzchołków (i dodatkowo odczytuje z tego samego pliku definicje krawędzi jako tekst). Po wczytaniu graf jest automatycznie wyświetlany w obszarze roboczym.

  \item \textbf{Wyb\'or algorytmu} --- mo\.zliwo\'s\'c wyboru algorytmu obliczeniowego: Fruchterman–Reingold lub Tutte. Uruchomienie nast\k{e}puje przez przycisk w panelu narz\k{e}dziowym.

  \item \textbf{Modyfikacja widoku} --- zmiana skali (zoom), przesuwanie obszaru roboczego (pan), automatyczne dopasowanie (Autofit) oraz pokazywanie lub ukrywanie etykiet wierzcho\l{}k\'ow i wag kraw\k{e}dzi. Panel narz\k{e}dziowy mo\.zna zwin\k{a}\'c przyciskiem \texttt{\guillemotleft/\guillemotright} przy lewej kraw\k{e}dzi panelu.

  \item \textbf{Edycja wsp\'o\l{}rz\k{e}dnych} --- interaktywna zmiana po\l{}o\.zenia wierzcho\l{}ka za pomoc\k{a} myszy (metoda drag-and-drop) bezpo\'srednio w obszarze roboczym lub poprzez wpisanie precyzyjnych warto\'sci w panelu narz\k{e}dziowym.

  \item \textbf{Automatyczne skalowanie} --- po wczytaniu danych lub obliczeniu uk\l{}adu, je\'sli wierzcho\l{}ki wykraczaj\k{a} poza wirtualny obszar roboczy ($5000 \times 5000$), aplikacja wy\'swietla okno dialogowe z pytaniem, czy automatycznie przeskalowa\'c wsp\'o\l{}rz\k{e}dne. U\.zytkownik mo\.ze zaakceptowa\'c (\textbf{Tak}) lub odrzuci\'c (\textbf{Nie}).

  \item \textbf{Zapis wynik\'ow} --- eksport obliczonych wsp\'o\l{}rz\k{e}dnych do pliku tekstowego lub binarnego, z mo\.zliwo\'sci\k{a} wyboru lokalizacji i nazwy pliku. Operacja wykonywana jest synchronicznie; po jej zako\'nczeniu pasek statusu wy\'swietla potwierdzenie.

  \item \textbf{Informowanie o b\l{}\k{e}dach} --- wy\'swietlanie modalnych okien dialogowych z opisem problemu (b\l{}\k{e}d krytyczny na czerwono, ostrze\.zenia na \.z\'o\l{}to), bez otwierania dodatkowych okien.
\end{itemize}

% =======================================================
\section{Opis funkcji i interfejsu}
% =======================================================

Aplikacja, zgodnie ze specyfikacj\k{a}, jest sterowana zdarzeniami w ramach interfejsu GUI zbudowanego w bibliotece Java Swing. Ca\l{}y program dzia\l{}a w jednym g\l{}\'{o}wnym oknie. \.Zadne z podstawowych akcji nie otwiera osobnego niezale\.znego okna --- w zamian stosowane s\k{a}: wbudowane okna dialogowe dla komunikat\'ow i plik\'ow, a tak\.ze okna wyboru pliku jako modalne nakładki.

\subsection{Schemat uk\l{}adu interfejsu}

Poni\.zszy diagram przedstawia uk\l{}ad komponent\'ow w g\l{}\'{o}wnym oknie aplikacji:

\begin{figure}[H]
\centering
\begin{tikzpicture}[
  font=\small\sffamily,
  box/.style={draw=primary, fill=lightbg, rounded corners=3pt,
              minimum width=#1, minimum height=0.7cm, align=center},
  box/.default=3cm,
  outerbox/.style={draw=rulegray, rounded corners=4pt, fill=white,
                   inner sep=8pt},
  lbl/.style={font=\scriptsize\ttfamily\color{rulegray}}
]

\node[outerbox, minimum width=13.5cm, minimum height=9.5cm] (win) at (0,0) {};
\node[lbl, above=2pt of win.north west, anchor=south west]
      {G\l{}\'owne okno aplikacji};

\node[box=13cm, fill=codebg, minimum height=0.6cm] (menu) at (0, 4.1)
      {\textbf{Pasek menu} \quad Plik ~|~ Pomoc};

\node[box=9cm, minimum height=6.8cm, fill=white, draw=accent]
      (canvas) at (-2.1, 0.2) {\textbf{Obszar roboczy}\\[4pt]renderowanie i podpowiedzi};

% Przycisk zwijania panelu narzędziowego
\node[box=0.4cm, minimum height=6.8cm, fill=codebg, font=\scriptsize]
      (toggle) at (2.95, 0.2) {\rotatebox{90}{\guillemotleft/\guillemotright}};

\node[box=3.5cm, minimum height=6.8cm, fill=lightbg]
      (tools) at (4.9, 0.2) {};
\node[font=\small\bfseries, above=0pt of tools.north, yshift=-18pt] {Panel Narz\k{e}dzi};

\node[box=3cm, minimum height=0.55cm, font=\scriptsize] (alg) at (4.9, 2.3)
      {Wyb\'or algorytmu};
\node[box=3cm, minimum height=0.55cm, font=\scriptsize] (run) at (4.9, 1.55)
      {Uruchom};
\node[box=3cm, minimum height=0.55cm, font=\scriptsize] (fit) at (4.9, 0.8)
      {Autofit};
\node[box=3cm, minimum height=0.55cm, font=\scriptsize] (zoom) at (4.9, 0.05)
      {Zoom +/-};
\node[box=3cm, minimum height=0.55cm, font=\scriptsize] (coord) at (4.9, -0.7)
      {Wsp. X / Y};
\node[box=3cm, minimum height=0.55cm, font=\scriptsize] (apply) at (4.9,-1.5)
      {Zastosuj};
\node[box=3cm, minimum height=0.55cm, font=\scriptsize] (vis) at (4.9,-2.3)
      {Etykiety / Wagi/ Pan};

\node[box=13cm, fill=codebg, minimum height=0.6cm, font=\scriptsize]
      (status) at (0, -3.7)
      {\textbf{Pasek statusu:} Stan aplikacji | Wierzchołki: <liczba> | Krawędzie: <liczba>};

\end{tikzpicture}
\caption{Uk\l{}ad komponent\'ow z uwzgl\k{e}dnieniem statystyk, funkcji Autofit i przycisku zwijania panelu.}
\end{figure}


\subsection{Pasek menu}

Opcje znajdujące się pod menu \texttt{Plik}:
\begin{itemize}
  \item \textbf{Wczytaj graf (Tekstowo)} --- otwiera modalne okno wyboru pliku nad głównym oknem. Użytkownik może wybrać plik tekstowy (\texttt{*.txt}) zawierający listę krawędzi. Po wyborze pliku i potwierdzeniu, aplikacja odczytuje definicje krawędzi (w formacie: nazwa, źródło, cel, waga). Wierzchołki, które nie zostały jawnie zdefiniowane, są tworzone automatycznie z domyślnymi współrzędnymi (0,0). Następnie graf jest renderowany w obszarze roboczym.
  
  \item \textbf{Wczytaj graf (Binarnie)} --- otwiera modalne okno wyboru pliku binarnego (\texttt{*.bin}). Aplikacja wczytuje współrzędne wierzchołków z binarnego nagłówka (format zgodny z wyjściem programu C), po czym z tego samego pliku próbuje odczytać listę krawędzi jako tekst. Dzięki temu możliwe jest wczytanie grafu, którego dane zostały wcześniej zapisane w formacie hybrydowym (binarne współrzędne + tekstowe krawędzie). Jeśli plik zawiera wyłącznie dane binarne, odczyt krawędzi może się nie powieść.

  \item \textbf{Zapisz wyniki (Tekstowo)} --- otwiera modalne okno zapisu pliku pozwalające użytkownikowi wybrać lokalizację i nazwę pliku tekstowego (\texttt{*.txt}). Zapisuje wyłącznie aktualne współrzędne wszystkich wierzchołków (ID, X, Y) w formacie tekstowym. Operacja jest synchroniczna; po jej zakończeniu pasek statusu wyświetla potwierdzenie.

  \item \textbf{Zapisz wyniki (Binarnie)} --- analogicznie do powyższego, ale z rozszerzeniem \texttt{*.bin}. Zapisuje współrzędne wierzchołków w formacie binarnym (little-endian, 20 bajtów na wierzchołek: int ID, double X, double Y). Format ten jest zgodny z oczekiwaniami silnika C.
\end{itemize}

\textbf{Uwaga:} Aplikacja nie oferuje funkcji "Zapisz graf" jako oddzielnej operacji. Dostępne są wyłącznie opcje eksportu samych współrzędnych (wyników obliczeń). Pełna definicja grafu (krawędzie wraz z wagami) pozostaje w oryginalnym pliku wejściowym i nie jest modyfikowana przez program.

\subsection{Panel narz\k{e}dziowy}

Panel ten daje pe\l{}n\k{a} kontrol\k{e} nad prezentowanymi danymi oraz umo\.zliwia uruchamianie algorytm\'ow przez program C:

\begin{itemize}
    \item \textbf{Wyb\'or algorytmu} --- rozwijana lista z opcjami \texttt{Fruchterman-Reingold} i \texttt{Tutte} (pe\l{}ne nazwy wy\'swietlane w interfejsie). Wyb\'or jest mapowany na odpowiedni argument \texttt{-a} do \texttt{graf\_planarny} (odpowiednio \texttt{fr} lub \texttt{tutte}). Klikni\k{e}cie przycisku \textbf{Uruchom} buduje i wykonuje polecenie, np.:
  \begin{cmdbox}
core/graf\_planarny -i wejscie.txt -o wynik.txt -a tutte -n 500
  \end{cmdbox}
  Po zako\'nczeniu procesu wynik jest automatycznie wczytywany i od\'swie\.zany w tym samym obszarze roboczym.

  \item \textbf{Liczba iteracji} --- pole tekstowe okre\'slaj\k{a}ce liczb\k{e} iteracji algorytmu Fruchtermana-Reingolda (domy\'slnie 100). W przypadku algorytmu Tutte pole jest automatycznie wy\l{}\k{a}czane i ustawiane na znak \texttt{-}, poniewa\.z algorytm ten nie korzysta z parametru iteracji.

  \item \textbf{Format danych (C)} --- rozwijana lista umo\.zliwiaj\k{a}ca wyb\'or formatu pliku tymczasowego przekazywanego do silnika C: \texttt{Tekstowy} lub \texttt{Binarny}. Opcja \texttt{Binarny} wymusza dodanie argumentu \texttt{-f bin} do wywo\l{}ania \texttt{graf\_planarny}.

  \item \textbf{Autofit} --- przycisk dopasowuj\k{a}cy automatycznie skal\k{e} i pozycj\k{e} widoku tak, aby ca\l{}y graf by\l{} widoczny w obszarze roboczym.

  \item \textbf{Modyfikacja wsp\'o\l{}rz\k{e}dnych} --- po klikni\k{e}ciu wierzcho\l{}ka na obszarze roboczym jego ID i wsp\'o\l{}rz\k{e}dne ($X$, $Y$) pojawiaj\k{a} si\k{e} w polach tekstowych w Panelu Narz\k{e}dzi. Przycisk \textbf{Zastosuj} przenosi wierzcho\l{}ek i od\'swie\.za widok.

  \item \textbf{Zmiana pola widzenia} --- suwak do zmiany skali (zoom in/out) oraz przycisk prze\l{}\k{a}czaj\k{a}cy tryb przesuwania (pan) obszaru roboczego.

  \item \textbf{Zwijanie panelu} --- po lewej kraw\k{e}dzi Panelu Narz\k{e}dzi znajduje si\k{e} przycisk ze strza\l{}k\k{a} (\texttt{\guillemotleft}/\texttt{\guillemotright}), kt\'ory pozwala schowa\'c lub ponownie wy\'swietli\'c ca\l{}y panel. Dzi\k{e}ki temu u\.zytkownik mo\.ze maksymalizowa\'c widoczno\'s\'c obszaru roboczego.

  \item \textbf{Interakcja z grafem} --- u\.zytkownik mo\.ze chwyci\'c dowolny wierzcho\l{}ek lewym przyciskiem myszy i przesun\k{e}\'c go w dowolne miejsce obszaru roboczego. Wsp\'o\l{}rz\k{e}dne w Panelu Narz\k{e}dzi s\k{a} aktualizowane w czasie rzeczywistym podczas przesuwania.

  \item \textbf{Spos\'ob wy\'swietlania} --- pola wyboru dla etykiet w\k{e}z\l{}\'ow i wag kraw\k{e}dzi.

  \item \textbf{Automatyczne skalowanie} --- po ka\.zdym wczytaniu pliku lub uruchomieniu algorytmu aplikacja sprawdza, czy wsp\'o\l{}rz\k{e}dne wierzcho\l{}k\'ow mieszcz\k{a} si\k{e} w obszarze roboczym ($12$--$4988$). Je\'sli kt\'orykolwiek wierzcho\l{}ek wykracza poza ten zakres, wy\'swietlane jest okno dialogowe z pytaniem: \textit{\guillemotleft{}Wykryto wierzcho\l{}ki poza obszarem roboczym. Czy chcesz je automatycznie przeskalowa\'c do obszaru $5000 \times 5000$?\guillemotright{}}. Po zatwierdzeniu (\textbf{Tak}) wsp\'o\l{}rz\k{e}dne s\k{a} normalizowane z zachowaniem proporcji.
\end{itemize}

% =======================================================
\section{Formaty wymienianych danych}
% =======================================================

Aplikacja rozróżnia dwa rodzaje plików: plik definicji grafu (lista krawędzi, używany jako wejście dla silnika C) oraz plik wynikowy (współrzędne wierzchołków, generowany przez silnik C lub przez funkcję zapisu aplikacji).

\subsection{Format pliku definicji grafu (lista krawędzi)}
Plik tekstowy (\texttt{*.txt}) zawierający wyłącznie definicje krawędzi – po jednej w każdym wierszu. Aplikacja akceptuje również komentarze (linie zaczynające się od \texttt{\#}) oraz puste linie, które są automatycznie usuwane przed przekazaniem do silnika C.

\textbf{Format pojedynczej linii z krawędzią:}

\begin{filebox}
<nazwa> <źródło> <cel> <waga>
\end{filebox}

\begin{itemize}
  \item \texttt{nazwa} --- identyfikator krawędzi (ciąg znaków bez spacji, np. \texttt{e1}).
  \item \texttt{źródło}, \texttt{cel} --- identyfikatory wierzchołków (liczby całkowite).
  \item \texttt{waga} --- liczba zmiennoprzecinkowa (z kropką dziesiętną).
\end{itemize}

Przykładowy plik wejściowy:
\begin{filebox}
\# Graf przykładowy (komentarz) \\
e1 \quad 1 \quad 2 \quad 1.5\\
e2 \quad 2 \quad 3 \quad 2.0\\
e3 \quad 3 \quad 1 \quad 0.8
\end{filebox}

\textbf{Uwaga:} Plik ten nie zawiera nagłówka ani osobnej listy wierzchołków – są one wyciągane dynamicznie z definicji krawędzi.

\subsection{Format pliku wynikowego (współrzędne wierzchołków)}

Plik generowany przez silnik C (lub przez funkcję zapisu aplikacji) zawiera wyłącznie współrzędne wierzchołków. Dostępne są dwie reprezentacje: tekstowa i binarna.

\subsubsection{Format tekstowy}

Pierwszy wiersz zawiera liczbę wierzchołków~$V$. Każdy kolejny wiersz reprezentuje jeden węzeł:

\begin{filebox}
V\\
<id> <x> <y>
\end{filebox}

\begin{itemize}
  \item \texttt{id} --- identyfikator węzła (liczba całkowita).
  \item \texttt{x, y} --- liczby zmiennoprzecinkowe, precyzja do 4 miejsc po przecinku.
\end{itemize}

Przykładowy plik wynikowy:
\begin{filebox}
3\\
1 \quad0.0000 \quad0.0000\\
2 \quad1.0000 \quad0.0000\\
3 \quad1.0000 \quad1.0000
\end{filebox}

\subsubsection{Format binarny}

Struktura pliku \texttt{.bin} (bez separatorów), \textbf{20 bajtów na wierzchołek}:

\begin{enumerate}
  \item \textbf{Nagłówek}: 4 bajty --- liczba wierzchołków~$V$ (typ \texttt{int}, little-endian).
  \item \textbf{Rekord wierzchołka} (powtarzany $V$ razy):
  \begin{itemize}
    \item 4 bajty: identyfikator węzła (typ \texttt{int}, w Javie odpowiada nieujemnym wartościom \texttt{unsigned int} z C),
    \item 8 bajtów: współrzędna $X$ (typ \texttt{double} IEEE~754),
    \item 8 bajtów: współrzędna $Y$ (typ \texttt{double} IEEE~754).
  \end{itemize}
\end{enumerate}

Kolejność bajtów jest zgodna z little-endian (natywna dla platform x86/x86-64), co jest bezpośrednio obsługiwane przez klasę \texttt{FileParser} w Javie poprzez mechanizm odwracania kolejności bajtów.

% =======================================================
\section{Sytuacje wyjątkowe}
% =======================================================

Wszystkie komunikaty błędów i ostrzeżeń wyświetlane są jako modalne okna dialogowe nad głównym oknem aplikacji. Nie otwierają one nowego niezależnego okna – są zawsze przyklejone do okna głównego. Po kliknięciu \textbf{OK} aplikacja wraca do stabilnego działania.

Stosowane są dwa rodzaje wizualne komunikatów:
\begin{itemize}
  \item \textbf{Czerwony popup (błąd krytyczny)} – tytuł okna: \texttt{Błąd}.
  \item \textbf{Żółty popup (ostrzeżenie)} – tytuł okna: \texttt{Ostrzeżenie}.
\end{itemize}

\textbf{Uwaga:} Wszystkie błędy pochodzące z silnika C wyświetlane są z przedrostkiem \texttt{"Błąd silnika: "}. Błędy wczytywania plików opatrzone są przedrostkiem \texttt{"Błąd wczytywania: "}, a błędy zapisu – \texttt{"Błąd zapisu: "}. Wyjątkiem jest przeciągnięcie wierzchołka poza widoczny obszar roboczy – wtedy nie jest wyświetlane okno dialogowe, a jedynie tymczasowy komunikat na pasku statusu (opisany na końcu sekcji).

\subsection*{Błędy krytyczne (czerwony popup)}

\begin{itemize}
  \item \textbf{Nie można uruchomić silnika obliczeniowego} – Wywołanie algorytmu, gdy plik \texttt{graf\_planarny} nie istnieje, nie ma praw do wykonania lub wystąpił błąd wejścia/wyjścia.\\
  Komunikat: \textit{„Błąd silnika: Nie odnaleziono lub nie można uruchomić pliku wykonywalnego graf\_planarny. Upewnij się, że plik istnieje i posiada prawo do wykonania (np. chmod +x).”}

  \item \textbf{Silnik C nie może otworzyć pliku wejściowego (kod 2)} – Program \texttt{graf\_planarny} zwrócił kod 2.\\
  Komunikat: \textit{„Błąd silnika: Silnik C nie może otworzyć pliku wejściowego. Sprawdź, czy plik istnieje i czy masz uprawnienia do odczytu.”}

  \item \textbf{Błąd formatu danych w pliku (kod 3)} – Program \texttt{graf\_planarny} zwrócił kod 3.\\
  Komunikat: \textit{„Błąd silnika: Błąd formatu danych w pliku (kod 3). Upewnij się, że plik krawędzi zawiera poprawne linie (np. 'e1 1 2 1.0') i nie posiada nagłówków ani komentarzy.”}

  \item \textbf{Nie można utworzyć pliku wynikowego przez program C (kod 4)} – Program \texttt{graf\_planarny} zwrócił kod 4.\\
  Komunikat: \textit{„Błąd silnika: Nie można utworzyć pliku wynikowego przez program C. Sprawdź uprawnienia katalogu docelowego.”}

  \item \textbf{Błąd parametrów (kod 5)} – Program \texttt{graf\_planarny} zwrócił kod 5 (nieznany algorytm).\\
  Komunikat: \textit{„Błąd silnika: Błąd parametrów (kod 5): Silnik nie rozpoznał nazwy algorytmu. Upewnij się, że przesyłasz 'fr' lub 'tutte'.”}

  \item \textbf{Błąd alokacji pamięci w silniku C (kod 7)} – Program \texttt{graf\_planarny} zwrócił kod 7.\\
  Komunikat: \textit{„Błąd silnika: Silnik obliczeniowy zgłosił błąd alokacji pamięci (kod 7). Graf może być zbyt duży. Spróbuj zmniejszyć liczbę wierzchołków.”}

  \item \textbf{Graf wejściowy jest niespójny (kod 8)} – Program \texttt{graf\_planarny} zwrócił kod 8.\\
  Komunikat: \textit{„Błąd silnika: Graf wejściowy jest niespójny (zawiera izolowane składowe). Program C wymaga grafu spójnego. Popraw plik wejściowy.”}

  \item \textbf{Błąd numeryczny w symulacji (kod 9)} – Program \texttt{graf\_planarny} zwrócił kod 9.\\
  Komunikat: \textit{„Błąd silnika: Wystąpił błąd numeryczny w symulacji (kod 9). Spróbuj uruchomić algorytm ponownie lub zwiększyć liczbę iteracji.”}

  \item \textbf{Pętla własna w pliku wejściowym (kod 12)} – Program \texttt{graf\_planarny} zwrócił kod 12.\\
  Komunikat: \textit{„Błąd silnika: Plik wejściowy zawiera pętlę własną (krawędź do samego siebie). Algorytmy nie obsługują tego rodzaju połączeń.”}

  \item \textbf{Graf zbyt gęsty / nieplanarny (kod 13)} – Program \texttt{graf\_planarny} zwrócił kod 13.\\
  Komunikat: \textit{„Błąd silnika: Graf jest zbyt gęsty, by mógł być planarny (E > 3V - 6). Zmniejsz liczbę krawędzi lub użyj innego grafu.”}

  \item \textbf{Nieoczekiwany błąd silnika} – Program \texttt{graf\_planarny} zwrócił kod inny niż wymienione powyżej.\\
  Komunikat: \textit{„Błąd silnika: Nieoczekiwany błąd silnika obliczeniowego (kod X).”} (gdzie X to kod błędu)

  \item \textbf{Błąd wczytywania pliku (nieprawidłowy format)} – Wystąpił błąd podczas parsowania pliku tekstowego (przy wczytywaniu grafu przez menu lub przez silnik C). Komunikat pochodzi z klasy \texttt{FileParser}.\\
  Komunikat: \textit{„Błąd wczytywania pliku! Oczekiwano danych w formacie: <węzeł> <x> <y>. Odnaleziono niedozwolone znaki w linii N.”} (gdzie N to numer linii)

  \item \textbf{Błąd wczytywania (inny)} – Inny wyjątek podczas wczytywania pliku, np. brak pliku.\\
  Komunikat: \textit{„Błąd wczytywania: ”} + szczegółowa treść błędu systemowego.

  \item \textbf{Plik binarny uszkodzony lub niekompletny} – Wczytywanie pliku \texttt{.bin}, gdy jego długość nie zgadza się z deklarowaną liczbą wierzchołków.\\
  Komunikat: \textit{„Błąd wczytywania: Plik binarny jest uszkodzony lub niekompletny. Oczekiwano N bajtów, znaleziono M bajtów.”}

  \item \textbf{Błąd zapisu pliku} – Wystąpił błąd wejścia/wyjścia podczas zapisywania współrzędnych.\\
  Komunikat: \textit{„Błąd zapisu: ”} + szczegółowa treść błędu systemowego.
\end{itemize}

\subsection*{Ostrzeżenia (żółty popup)}

\begin{itemize}
  \item \textbf{Brak pliku wejściowego przed uruchomieniem algorytmu} – Kliknięcie \textbf{Uruchom}, gdy użytkownik nie wybrał jeszcze żadnego pliku z listą krawędzi.\\
  Komunikat: \textit{„Najpierw wczytaj plik z listą krawędzi!”}

  \item \textbf{Graf jest pusty} – Próba uruchomienia algorytmu na grafie bez wierzchołków.\\
  Komunikat: \textit{„Graf jest pusty. Wczytaj poprawne dane przed uruchomieniem algorytmu.”}

  \item \textbf{Nieprawidłowa liczba iteracji} – W polu iteracji wpisano wartość nie będącą dodatnią liczbą całkowitą.\\
  Komunikat: \textit{„Liczba iteracji musi być całkowitą liczbą dodatnią!”}

  \item \textbf{Współrzędna poza dopuszczalnym zakresem} – Kliknięcie \textbf{Zastosuj}, gdy wprowadzona współrzędna jest mniejsza niż 12 lub większa niż 4988.\\
  Komunikat: \textit{„Współrzędne muszą znajdować się w obszarze roboczym (12 - 4988).”}

  \item \textbf{Niepoprawna wartość współrzędnej (pole tekstowe)} – Kliknięcie \textbf{Zastosuj}, gdy w polach X lub Y znajduje się wartość nieskończona lub \texttt{NaN}.\\
  Komunikat: \textit{„Współrzędne muszą być skończonymi liczbami rzeczywistymi.”}

  \item \textbf{Niepoprawny format liczby w polu współrzędnych} – Gdy wpisana wartość nie jest liczbą.\\
  Komunikat: \textit{„Wprowadzona wartość nie jest poprawną liczbą. Użyj formatu np. 123.45”}
\end{itemize}

\subsection*{Komunikat na pasku statusu (bez okna dialogowego)}

\textbf{Próba umieszczenia wierzchołka poza widocznym obszarem roboczym} – Przeciągnięcie i upuszczenie wierzchołka w miejsce poza obszarem roboczym. Wierzchołek zostaje zablokowany na najbliższym brzegu. Na pasku statusu pojawia się na 3 sekundy komunikat: \\
\textit{„Uwaga: Próba przesunięcia wierzchołka poza obszar roboczy!”} \\
Po tym czasie pasek statusu wraca do standardowego podsumowania (liczba wierzchołków i krawędzi).

Aby zminimalizować występowanie błędów po stronie interfejsu, komponenty są blokowane z zachowaniem spójności --- opcje eksportu, \textbf{Uruchom} oraz \textbf{Zastosuj} są dezaktywowane, dopóki nie zostanie wczytany model grafu.


% ============================================================
% ============================================================
%   CZĘŚĆ II: DOKUMENTACJA IMPLEMENTACYJNA
% ============================================================
% ============================================================

\partpage{Część II}{Dokumentacja implementacyjna}

% ============================================================
\section{Wstęp do dokumentacji implementacyjnej}
% ============================================================

Dokument stanowi przewodnik po kodzie źródłowym aplikacji okienkowej
\texttt{Wizualizacja grafu planarnego} --- narzędzia napisanego w języku
Java (biblioteka Swing), pełniącego rolę graficznej nakładki na program
obliczeniowy wykonany w języku C. Dokumentacja opisuje podział systemu
na moduły, wewnętrzne struktury danych, architekturę interfejsu
użytkownika, przepływ sterowania, zasady komunikacji między komponentami
oraz zastosowane wzorce projektowe. Jej celem jest ułatwienie zrozumienia,
pielęgnacji i rozbudowy kodu.

Aplikacja Java nie implementuje własnych algorytmów obliczeniowych.
Cała logika matematyczna pozostaje po stronie zewnętrznego
programu \texttt{graf\_planarny}, uruchamianego jako osobny proces
przez klasę \texttt{CalculationEngine}. Aplikacja odpowiada wyłącznie
za prezentację grafu, interakcję z użytkownikiem oraz przekazywanie
danych do i z silnika C.

% ============================================================
\section{Struktura plików i podział na moduły}
% ============================================================

Kod źródłowy podzielony jest na dziesięć modułów, każdy reprezentowany
przez osobny plik \texttt{.java}. Jedynymi typami danych współdzielonymi
między modułami są \texttt{Graph}, \texttt{Vertex} i \texttt{Edge},
reprezentujące model grafu.

\begin{center}
\begin{tabular}{>{\ttfamily}l p{9.2cm}}
\toprule
\normalfont\textbf{Plik} & \textbf{Odpowiedzialno\'s\'c} \\
\midrule
Main.java            & Punkt wejścia aplikacji. Główne okno \texttt{JFrame}, koordynacja wszystkich komponentów GUI, obsługa zdarzeń wysokiego poziomu (wczytywanie, zapis, uruchamianie algorytmu). \\
\addlinespace
Graph.java           & Definicja struktury grafu: mapa wierzchołków \texttt{TreeMap<Integer, Vertex>}, lista krawędzi \texttt{List<Edge>}, obliczanie obwiedni (bounds). \\
\addlinespace
Vertex.java          & Pojedynczy wierzchołek: identyfikator \texttt{id} oraz współrzędne \texttt{x}, \texttt{y} z akcesorami. \\
\addlinespace
Edge.java            & Pojedyncza krawędź: nazwa, referencje do wierzchołków źródłowego i docelowego, waga. \\
\addlinespace
FileParser.java      & Operacje wejścia/wyjścia: wczytywanie grafu (tekstowe i binarne), zapis wyników (tekstowy i binarny), konwersja little-endian $\leftrightarrow$ big-endian. \\
\addlinespace
CalculationEngine.java & Uruchamianie zewnętrznego programu \texttt{graf\_planarny} (C) jako procesu potomnego, interpretacja kodów powrotu. \\
\addlinespace
GraphPanel.java      & Komponent \texttt{JPanel} odpowiedzialny za renderowanie grafu, obsługę myszy (zaznaczanie, przeciąganie wierzchołków, tryb przesuwania), skalowanie i autofit. \\
\addlinespace
AppToolPanel.java    & Panel narzędziowy (prawy): wybór algorytmu, liczba iteracji, format danych, zoom, edycja współrzędnych, widoczność etykiet i wag. \\
\addlinespace
AppMenuBar.java      & Pasek menu: opcje wczytywania grafu (tekstowo i binarnie), zapisu wyników (tekstowo/binarnie) oraz okno pomocy. \\
\addlinespace
AppStatusPanel.java  & Pasek statusu (dolny): wyświetlanie bieżącego stanu aplikacji, liczby wierzchołków i krawędzi oraz tymczasowych komunikatów o błędach. \\
\bottomrule
\end{tabular}
\end{center}

\subsection{Zależności między modułami}

Klasa \texttt{Main} zna wszystkie pozostałe komponenty i~zarządza
komunikacją między nimi. Pozostałe powiązania:

\begin{itemize}
  \item \texttt{GraphPanel} korzysta z~\texttt{Graph} (dane grafu)
    i~\texttt{Main} (powiadamianie o~zdarzeniach),
  \item \texttt{AppToolPanel} korzysta z~\texttt{Main}
    i~\texttt{GraphPanel} (sterowanie zoomem),
  \item \texttt{AppMenuBar} i~\texttt{AppStatusPanel} zgłaszają
    akcje do \texttt{Main} przez listenery,
  \item \texttt{FileParser}, \texttt{CalculationEngine} i~\texttt{Graph}
    to klasy pomocnicze --- nie wiedzą o~GUI,
  \item \texttt{Vertex} i~\texttt{Edge} to proste struktury danych
    używane przez \texttt{Graph}, \texttt{GraphPanel}
    i~\texttt{FileParser}.
\end{itemize}

% ============================================================
\section{Struktury danych}
% ============================================================

Podstawowe struktury danych reprezentujące graf zdefiniowane są
w osobnych plikach \texttt{Vertex.java}, \texttt{Edge.java}
oraz \texttt{Graph.java}. W odróżnieniu od części C, gdzie dominują
tablice dynamiczne i ręczne zarządzanie pamięcią, w części Java
wykorzystywane są kolekcje z biblioteki standardowej
(\texttt{TreeMap}, \texttt{ArrayList}).

\subsection{Vertex --- wierzchołek grafu}

\begin{filebox}
public class Vertex \{\\
\quad int id;\\
\quad double x, y;\\
\\
\quad public Vertex(int id, double x, double y);\\
\quad public double getX();\\
\quad public double getY();\\
\quad public void setX(double x);\\
\quad public void setY(double y);\\
\}
\end{filebox}

Klasa przechowuje identyfikator wierzchołka (\texttt{id}) oraz jego
bieżące współrzędne $(x, y)$ w wirtualnym obszarze roboczym
$[0, 5000] \times [0, 5000]$. Pola są pakietowo dostępne (modyfikator
domyślny), co umożliwia szybki dostęp z \texttt{GraphPanel} (renderowanie)
bez narzutu wywołań getterów/setterów w pętli rysowania. Metody
\texttt{getX()} i \texttt{getY()} są zachowane dla klas zewnętrznych
względem pakietu.

\subsection{Edge --- krawędź grafu}

\begin{filebox}
public class Edge \{\\
\quad private String name;\\
\quad private Vertex source;\\
\quad private Vertex target;\\
\quad private double weight;\\
\\
\quad public Edge(String name, Vertex source,\\
\quad \qquad\qquad\; Vertex target, double weight);\\
\quad public String getName();\\
\quad public double getWeight();\\
\quad public Vertex getSource();\\
\quad public Vertex getTarget();\\
\}
\end{filebox}

Klasa przechowuje nazwę krawędzi, referencje do obiektów \texttt{Vertex}
(źródło i cel) oraz wagę (liczba zmiennoprzecinkowa podwójnej precyzji).
Referencje do wierzchołków umożliwiają bezpośredni dostęp do współrzędnych
podczas rysowania, bez konieczności wyszukiwania w mapie.

\subsection{Graph --- główna struktura grafu}

\begin{filebox}
public class Graph \{\\
\quad Map<Integer, Vertex> vertices = new TreeMap<>();\\
\quad List<Edge> edges = new ArrayList<>();\\
\quad private double minX, maxX, minY, maxY;\\
\\
\quad public void addVertex(int id, double x, double y);\\
\quad public void addEdge(String name, int sId, int tId,\\
\quad \qquad\qquad\qquad\quad\;\; double weight);\\
\quad public int getVertexCount();\\
\quad public int getEdgesCount();\\
\quad public Map<Integer, Vertex> getVertices();\\
\quad public void calculateBounds();\\
\quad public double getMinX();\\
\quad public double getMaxX();\\
\quad public double getMinY();\\
\quad public double getMaxY();\\
\quad public void clear();\\
\}
\end{filebox}

Wierzchołki przechowywane są w mapie \texttt{TreeMap<Integer, Vertex>},
co gwarantuje kolejność rosnącą według identyfikatora podczas iteracji
(np. przy zapisie wyników) i zapewnia średni czas dostępu $O(\log V)$.
Krawędzie przechowywane są w liście \texttt{ArrayList<Edge>}.

Metoda \texttt{addEdge} automatycznie wyszukuje wierzchołki źródłowy
i docelowy w mapie; jeśli oba istnieją, tworzy nowy obiekt \texttt{Edge}
i dodaje go do listy. W przeciwnym razie krawędź jest ignorowana
(co zapobiega niespójności danych).

Pola \texttt{minX}, \texttt{maxX}, \texttt{minY}, \texttt{maxY}
przechowują obwiednię grafu obliczaną przez \texttt{calculateBounds()}.
Są one wykorzystywane przez \texttt{GraphPanel} do automatycznego
skalowania (\texttt{autofit()}).

% ============================================================
\section{Moduł \texttt{FileParser} --- wejście/wyjście}
% ============================================================

Klasa \texttt{FileParser} zawiera całą logikę związaną z odczytem
i zapisem plików. Wszystkie publiczne metody operują na przekazanym
obiekcie \texttt{Graph}, modyfikując go przez jego metody publiczne.

\subsection{Metody prywatne --- wczytywanie}

Wczytywanie danych odbywa się przez trzy metody prywatne, wywoływane
w zależności od formatu pliku:

\begin{itemize}
  \item \texttt{readVertices(String path, Graph graph)} --- wczytuje
    plik tekstowy zawierający listę wierzchołków (pierwsza linia: liczba
    wierzchołków; kolejne: \texttt{id x y}). Używana do odczytu plików
    wynikowych generowanych przez silnik C lub przez funkcję zapisu.
    Format wejściowy używa kropki dziesiętnej (\texttt{Locale.US}).

  \item \texttt{readEdges(String path, Graph graph)} --- wczytuje
    plik tekstowy z listą krawędzi w formacie \texttt{nazwa zrodlo cel waga}.
    Jeśli napotkany wierzchołek nie istnieje jeszcze w grafie, jest
    automatycznie tworzony z domyślną pozycją $(0, 0)$. Metoda pomija
    linie niepasujące do formatu (zabezpieczenie \texttt{hasNextInt()}
    przed wyjątkiem).

  \item \texttt{readBinary(String path, Graph graph)} --- wczytuje
    binarny plik wynikowy w formacie little-endian (zgodnym z wyjściem
    programu C na platformie x86). Struktura pliku: 4 bajty --- liczba
    wierzchołków $V$ (\texttt{int}), następnie $V$ rekordów po 20 bajtów
    każdy (\texttt{int id}, \texttt{double x}, \texttt{double y}).
    Konwersja little-endian $\rightarrow$ big-endian realizowana jest
    przez wewnętrzne odwracanie kolejności bajtów w buforze.
    Po wczytaniu wywoływane jest \texttt{calculateBounds()}.
\end{itemize}

\subsection{Metody publiczne}

\subsubsection{loadFullGraph}

\begin{cmdbox}
public void loadFullGraph(String nodePath, String edgePath,\\
\quad\qquad\qquad\qquad\quad\; Graph graph, boolean binary)\\
\quad\qquad\qquad\qquad\quad\; throws Exception;
\end{cmdbox}

\textbf{Parametry:}
\begin{itemize}
  \item \texttt{nodePath} --- ścieżka do pliku z listą wierzchołków
    (wynikowego z silnika C); dla trybu binarnego traktowana również
    jako źródło danych binarnych,
  \item \texttt{edgePath} --- ścieżka do pliku z listą krawędzi,
  \item \texttt{graph} --- obiekt grafu do wypełnienia,
  \item \texttt{binary} --- jeśli \texttt{true}, wierzchołki wczytywane
    są z pliku binarnego (\texttt{readBinary}); jeśli \texttt{false},
    z tekstowego (\texttt{readVertices}).
\end{itemize}

\textbf{Działanie:} Wczytuje wierzchołki (z formatu tekstowego lub
binarnego), a następnie krawędzie (zawsze z pliku tekstowego
\texttt{edgePath}). Czyszczenie grafu (\texttt{graph.clear()}) odbywa się wewnątrz metod \texttt{readVertices()} i \texttt{readBinary()}, nie bezpośrednio w \texttt{loadFullGraph()}.

\textbf{Zwraca:} nic; w przypadku błędu rzuca \texttt{Exception}
z opisem problemu.

\subsubsection{saveToText}

\begin{cmdbox}
public void saveToText(String path, Graph graph) throws Exception;
\end{cmdbox}

Zapisuje współrzędne wszystkich wierzchołków do pliku tekstowego.
Pierwsza linia zawiera liczbę wierzchołków. Każda kolejna linia:
\texttt{id x y} z precyzją 4 miejsc po przecinku
(\texttt{printf("\%d \%.4f \%.4f\textbackslash n")}). Używa \texttt{Locale.US}
dla wymuszenia kropki dziesiętnej.

\subsubsection{saveToBinary}

\begin{cmdbox}
public void saveToBinary(String path, Graph graph) throws Exception;
\end{cmdbox}

Zapisuje współrzędne do pliku binarnego w formacie little-endian
(kompatybilnym z programem C). Używa \texttt{ByteBuffer} z kolejnością
bajtów \texttt{LITTLE\_ENDIAN}. Struktura: 4 bajty --- liczba
wierzchołków, następnie rekordy 20-bajtowe (\texttt{int id},
\texttt{double x}, \texttt{double y}).

% ============================================================
\section{Moduł \texttt{CalculationEngine} --- silnik obliczeniowy}
% ============================================================

Klasa \texttt{CalculationEngine} odpowiada za uruchomienie zewnętrznego
programu \texttt{graf\_planarny} (C) jako procesu potomnego oraz za
interpretację jego kodu powrotu.

\subsection{runGraphAlgorithm}

\begin{cmdbox}
public void runGraphAlgorithm(String inputPath,\\
\quad\qquad\qquad\qquad\qquad\; String algorithm,\\
\quad\qquad\qquad\qquad\qquad\; int iterations,\\
\quad\qquad\qquad\qquad\qquad\; boolean binary,\\
\quad\qquad\qquad\qquad\qquad\; Graph graph)\\
\quad\qquad\qquad\qquad\qquad\; throws Exception;
\end{cmdbox}

\textbf{Parametry:}
\begin{itemize}
  \item \texttt{inputPath} --- ścieżka do pliku wejściowego z listą krawędzi,
  \item \texttt{algorithm} --- nazwa algorytmu (\texttt{"fr"} lub
    \texttt{"tutte"}); przekazywana jako argument \texttt{-a},
  \item \texttt{iterations} --- liczba iteracji (dla FR); przekazywana
    jako argument \texttt{-n},
  \item \texttt{binary} --- jeśli \texttt{true}, dodaje flagę
    \texttt{-f bin} do polecenia,
  \item \texttt{graph} --- obiekt grafu, do którego zostaną wczytane
    wyniki po pomyślnym zakończeniu procesu.
\end{itemize}

\textbf{Działanie:}
\begin{enumerate}
  \item Ustala ścieżkę do pliku wykonywalnego: \texttt{core/graf\_planarny}
    (Linux/macOS) lub \texttt{core/graf\_planarny.exe} (Windows).
  \item Buduje listę argumentów: \texttt{-i <inputPath>},
    \texttt{-o wynik.txt}, \texttt{-a <algorithm>},
    \texttt{-n <iterations>} oraz opcjonalnie \texttt{-f bin}.
  \item Uruchamia proces przez \texttt{ProcessBuilder}.
  \item Oczekuje na zakończenie procesu (\texttt{waitFor()}).
  \item Jeśli kod powrotu to 0, wywołuje \texttt{FileParser.loadFullGraph}
    dla pliku \texttt{wynik.txt} i ścieżki wejściowej.
  \item Jeśli kod powrotu $\neq$ 0, rzuca \texttt{Exception}
    z komunikatem odpowiadającym kodowi błędu (mapowanie kodów
    realizowane jest wewnątrz klasy).
  \item Jeśli proces nie może zostać uruchomiony (\texttt{IOException}),
    rzuca \texttt{Exception} z komunikatem o braku pliku wykonywalnego.
\end{enumerate}

% ============================================================
\section{Moduł \texttt{GraphPanel} --- panel renderowania grafu}
% ============================================================

Klasa \texttt{GraphPanel} rozszerza \texttt{JPanel} i odpowiada za całą
wizualną reprezentację grafu oraz interakcję myszą.

\subsection{Obszar roboczy i układ współrzędnych}

Wirtualny obszar roboczy ma stały rozmiar $5000 \times 5000$ pikseli
(\texttt{WORK\_WIDTH}, \texttt{WORK\_HEIGHT}). Współrzędne wierzchołków
przechowywane w obiektach \texttt{Vertex} odnoszą się do tego wirtualnego
obszaru. Podczas renderowania stosowane jest przekształcenie:

\begin{enumerate}
  \item \textbf{Translacja} --- środek wirtualnego obszaru roboczego
    jest wyrównywany ze środkiem widocznego panelu. Dodatkowo stosowane
    jest przesunięcie \texttt{offsetX}/\texttt{offsetY} (tryb pan).
  \item \textbf{Skalowanie} --- współczynnik \texttt{zoom} (domyślnie
    $0.5$, zakres od $\max(W_{\text{panelu}}/5000,\\\; H_{\text{panelu}}/5000)$
    do $2.0$) jest kontrolowany suwakiem w panelu narzędziowym.
\end{enumerate}

Transformacja realizowana jest przez wywołania \texttt{g2.translate(tx, ty)}
i \texttt{g2.scale(zoom, zoom)} przed pętlą rysowania. Dzięki temu
rysowanie wierzchołków i krawędzi odbywa się bezpośrednio na podstawie
współrzędnych wirtualnych.

\subsection{Obsługa myszy}

Wewnętrzna klasa anonimowa \texttt{MouseAdapter} obsługuje trzy rodzaje
interakcji:

\begin{itemize}
  \item \textbf{Kliknięcie (mousePressed)} --- jeśli tryb \texttt{panMode}
    jest aktywny, zapamiętuje pozycję myszy do późniejszego przesuwania.
    W przeciwnym razie przeszukuje wierzchołki w promieniu
    $2 \cdot V\_RAD$ (metoda \texttt{findVertex}) i ustawia znaleziony
    obiekt jako \texttt{selectedVertex}.

  \item \textbf{Przeciąganie (mouseDragged)} --- w trybie pan:
    aktualizuje \texttt{offsetX}/\texttt{offsetY} o wektor przesunięcia
    myszy i ogranicza przesunięcie do granic obszaru roboczego
    (\texttt{clampOffsets()}). W trybie edycji: przesuwa zaznaczony
    wierzchołek do pozycji kursora (z ograniczeniem do widocznego
    obszaru i do granic wirtualnego obszaru roboczego). Przekroczenie
    granic zgłaszane jest do \texttt{Main.notifyVertexOutOfBounds()}.
\end{itemize}

\subsection{Autofit}

Metoda \texttt{autofit()} oblicza obwiednię grafu, wyznacza skalę
tak, aby graf zajmował ok. 80\% widocznego obszaru, i odpowiednio
przeskalowuje współrzędne wszystkich wierzchołków. Po przeskalowaniu
graf jest centrowany w wirtualnym obszarze $5000 \times 5000$.

\subsection{Normalizacja do obszaru roboczego}

Metoda \texttt{normalizeToWorkspace()} skaluje i przesuwa współrzędne
wszystkich wierzchołków tak, aby zmieściły się w obszarze
$[12, 4988] \times [12, 4988]$ z zachowaniem proporcji. Wywoływana
jest warunkowo z metody \texttt{Main.updateUIState()} --- jeśli po wczytaniu pliku lub uruchomieniu algorytmu wierzchołki wykraczają poza granice obszaru roboczego, wyświetlany jest dialog potwierdzenia (\texttt{JOptionPane.YES\_NO\_OPTION}), a po zatwierdzeniu przez użytkownika wywoływana jest \texttt{normalizeToWorkspace()}.

\subsection{Renderowanie}

Metoda \texttt{paintComponent} wykonuje kolejno:
\begin{enumerate}
  \item Ograniczenie przesunięcia kamery (\texttt{clampOffsets()}).
  \item Obliczenie translacji. Zastosowanie \texttt{translate} i \texttt{scale}.
  \item Narysowanie tła obszaru roboczego (jasnoszary prostokąt
    $5000 \times 5000$) i jego ramki.
  \item Narysowanie wszystkich krawędzi jako linii (ciemnoszary).
    Opcjonalnie wyświetlenie wag w punkcie środkowym krawędzi
    (gdy \texttt{showWeights == true}).
  \item Narysowanie wszystkich wierzchołków jako okręgów o promieniu
    \texttt{V\_RAD = 12}. Wierzchołek zaznaczony (\texttt{selectedVertex})
    wypełniany jest kolorem białym, pozostałe --- czarnym.
    Opcjonalnie wyświetlenie etykiet (gdy \texttt{showLabels == true}).
\end{enumerate}

% ============================================================
\section{Pozostałe moduły GUI}
% ============================================================

\subsection{AppMenuBar}

Rozszerza \texttt{JMenuBar}. Tworzy dwa menu:

\begin{itemize}
  \item \textbf{Plik} --- zawiera opcje:
    \begin{itemize}
      \item \texttt{Wczytaj graf (Tekstowo)} --- wywołuje \texttt{parent.openFile(false)},
      \item \texttt{Wczytaj graf (Binarnie)} --- wywołuje \texttt{parent.openFile(true)},
      \item separator,
      \item \texttt{Zapisz wyniki (Tekstowo)} --- wywołuje
        \texttt{parent.saveAction(false)},
      \item \texttt{Zapisz wyniki (Binarnie)} --- wywołuje
        \texttt{parent.saveAction(true)}.
    \end{itemize}
  \item \textbf{Pomoc} --- zawiera opcję \texttt{Instrukcja i dokumentacja},
    wyświetlającą modalne okno dialogowe ze skróconą instrukcją obsługi
    i odesłaniem do pełnej dokumentacji.
\end{itemize}

\subsection{AppToolPanel}

Rozszerza \texttt{JPanel} z układem \texttt{BoxLayout} (oś Y).
Podzielony na cztery sekcje wizualne:

\begin{enumerate}
  \item \textbf{Ustawienia silnika} --- rozwijana lista algorytmów
    (\texttt{Fruchterman-Reingold} / \texttt{Tutte}), pole iteracji
    (nieaktywne dla Tutte), rozwijana lista formatu danych
    (\texttt{Tekstowy} / \texttt{Binarny}), przycisk \texttt{Uruchom}.
  \item \textbf{Widok} --- suwak zoom (zakres $5\%$--$200\%$ na sliderze,
    efektywne minimum ograniczane dynamicznie przez \texttt{GraphPanel.getMinZoom()},
    które zapobiega oddaleniu poniżej rozmiaru widocznego panelu),
    przycisk \texttt{Autofit}.
  \item \textbf{Edycja wierzchołka} --- pola tekstowe $X$ i $Y$,
    etykieta z ID zaznaczonego wierzchołka, przycisk \texttt{Zastosuj}.
  \item \textbf{Widoczność} --- pola wyboru: \texttt{Etykiety},
    \texttt{Wagi}, \texttt{Przesuwanie (Pan)}.
\end{enumerate}

Komunikacja z \texttt{Main} odbywa się przez referencję \texttt{parent},
a z \texttt{GraphPanel} przez referencję \texttt{graphPanel} (do
bezpośredniego ustawiania zoomu). Metody dostępowe
(\texttt{getSelectedAlgorithm()}, \texttt{isBinaryFormat()},
\texttt{getIterations()}, \texttt{setIterationsEnabled()},
\texttt{setVertexInfo()},
\texttt{getXText()}, \texttt{getYText()},
\texttt{setXText()}, \texttt{setYText()},
\texttt{updateZoomSlider()}) umożliwiają \texttt{Main}
pobieranie i ustawianie wartości kontrolek.

\subsection{AppStatusPanel}

Rozszerza \texttt{JPanel}. Wyświetla bieżący stan aplikacji w formacie:

\begin{filebox}
Stan: <komunikat> | Wierzchołki: <N> | Krawędzie: <M>
\end{filebox}

Metoda \texttt{notifyError} wyświetla tymczasowy komunikat o błędzie
(kolor czerwony), który po 3 sekundach (\texttt{Timer}) automatycznie
powraca do poprzedniego stanu. Mechanizm ten stosowany jest wyłącznie
dla ostrzeżeń o próbie przesunięcia wierzchołka poza obszar roboczy ---
wszystkie pozostałe błędy obsługiwane są przez modalne okna dialogowe
\texttt{JOptionPane} w klasie \texttt{Main}.

% ============================================================
\section{Moduł \texttt{Main} --- punkt wejścia i przepływ sterowania}
% ============================================================

Klasa \texttt{Main} rozszerza \texttt{JFrame} i pełni rolę:
\begin{itemize}
  \item punktu wejścia aplikacji (metoda \texttt{main}),
  \item głównego okna (tytuł: \texttt{"Wizualizacja grafu planarnego"}),
  \item mediatora między wszystkimi komponentami GUI,
  \item kontrolera przepływu sterowania.
\end{itemize}

\subsection{Metoda \texttt{main}}

\begin{cmdbox}
public static void main(String[] args);
\end{cmdbox}

Wywołuje \texttt{SwingUtilities.invokeLater()} z konstruktorem
\texttt{Main}, co zapewnia utworzenie GUI w wątku EDT (Event Dispatch
Thread), zgodnie z wymaganiami Swing.

\subsection{Konstruktor \texttt{Main()}}

Wykonuje inicjalizację wszystkich komponentów:
\begin{enumerate}
  \item Tworzy obiekty modelu (\texttt{Graph}), parsera (\texttt{FileParser})
    i silnika (\texttt{CalculationEngine}).
  \item Tworzy \texttt{GraphPanel} i przekazuje mu referencje do grafu
    oraz do siebie.
  \item Ustawia \texttt{AppMenuBar} jako pasek menu okna.
  \item Tworzy \texttt{AppToolPanel} (prawy panel narzędziowy)
    i \texttt{AppStatusPanel} (dolny pasek statusu).
  \item Rozmieszcza komponenty w układzie \texttt{BorderLayout}:
    \texttt{GraphPanel} w centrum, kontener \texttt{rightContainer}
    (zawierający przycisk zwijania \texttt{toggleBtn} i \texttt{AppToolPanel})
    po prawej stronie, pasek statusu na dole. Przycisk \texttt{toggleBtn}
    (symbole \texttt{\guillemotleft}/\texttt{\guillemotright}) przełącza widoczność panelu narzędziowego.
  \item Dezaktywuje panel narzędziowy do czasu wczytania pierwszego pliku.
  \item Wywołuje \texttt{graphPanel.autofit()} po wyrenderowaniu UI.
\end{enumerate}

\subsection{Główne metody sterujące}

\subsubsection{openFile}

\begin{cmdbox}
public void openFile(boolean binary);
\end{cmdbox}

Otwiera modalne okno wyboru pliku (\texttt{JFileChooser}). Po wyborze:
\begin{enumerate}
  \item Czyści graf (\texttt{graph.clear()}).
  \item W trybie binarnym wczytuje dane przez
    \texttt{parser.loadFullGraph(..., true)}.
  \item W trybie tekstowym:
    \begin{enumerate}
      \item Wywołuje prywatną metodę pomocniczą \texttt{prepareCleanedFile()},
        która tworzy plik tymczasowy pozbawiony komentarzy (linii zaczynających się
        od \texttt{\#}) i pustych linii. Plik tymczasowy oznaczany jest przez
        \texttt{deleteOnExit()}.
      \item Wywołuje przez refleksję prywatną metodę \texttt{readEdges}
        klasy \texttt{FileParser} dla oczyszczonego pliku.
    \end{enumerate}
  \item Aktywuje panel narzędziowy.
  \item Aktualizuje stan UI (\texttt{updateUIState()}).
\end{enumerate}
W przypadku błędu wyświetla okno dialogowe z komunikatem.

\subsubsection{runAlgorithm}

\begin{cmdbox}
public void runAlgorithm();
\end{cmdbox}

Wywoływana przez przycisk \texttt{Uruchom} w panelu narzędziowym.
Waliduje stan (czy wczytano plik, czy graf nie jest pusty), parsuje
parametry (algorytm, liczba iteracji, format), a następnie wywołuje
\texttt{engine.runGraphAlgorithm()}. Po pomyślnym zakończeniu odświeża
widok (\texttt{autofit()}) i pasek statusu. W przypadku błędu wyświetla
komunikat z prefiksem \texttt{"Błąd silnika: "}.

\subsubsection{applyCoords}

\begin{cmdbox}
public void applyCoords();
\end{cmdbox}

Wywoływana przez przycisk \texttt{Zastosuj} w panelu edycji wierzchołka.
Waliduje wprowadzone współrzędne: sprawdza, czy mieszczą się w obszarze
roboczym $[12, 4988]$, a następnie czy są skończonymi liczbami
rzeczywistymi (tj.\ nie są \texttt{NaN} ani nieskończone).
Następnie aktualizuje pozycję zaznaczonego wierzchołka i odświeża widok.
W przypadku błędu formatu lub przekroczenia zakresu wyświetla ostrzeżenie.

\subsubsection{saveAction}

\begin{cmdbox}
public void saveAction(boolean binary);
\end{cmdbox}

Otwiera modalne okno zapisu pliku (\texttt{JFileChooser}).
W zależności od flagi \texttt{binary} wywołuje \texttt{parser.saveToText()}
lub \texttt{parser.saveToBinary()}.

\subsubsection{updateUIState}

\begin{cmdbox}
private void updateUIState(String msg);
\end{cmdbox}

Wspólna metoda aktualizacji stanu całego interfejsu, wywoływana po każdej
operacji modyfikującej graf (wczytanie pliku, uruchomienie algorytmu).
Wykonuje kolejno:
\begin{enumerate}
  \item Aktualizuje pasek statusu (komunikat, liczba wierzchołków i krawędzi).
  \item Wywołuje \texttt{graphPanel.autofit()} w celu dopasowania widoku.
  \item Sprawdza, czy współrzędne wierzchołków mieszczą się w~obszarze roboczym ($12$--$4988$). Jeśli nie, wyświetla modalne okno dialogowe (\texttt{JOptionPane.YES\_NO\_OPTION}) z~pytaniem o~automatyczne przeskalowanie.
  \item Po zatwierdzeniu wywołuje \texttt{graphPanel.normalizeToWorkspace()}.
  \item Synchronizuje suwak zoomu z~bieżącą wartością.
\end{enumerate}

% ============================================================
\section{Zarządzanie pamięcią i zasobami}
% ============================================================

Zarządzanie pamięcią w części Java opiera się
na automatycznym odśmiecaniu (Garbage Collector). Niemniej program
przestrzega kilku zasad:

\begin{itemize}
  \item \textbf{Referencje} --- obiekty \texttt{Vertex} są współdzielone
    między mapą \texttt{Graph.vertices} a listą krawędzi
    (\texttt{Edge.source}, \texttt{Edge.target}). Dzięki temu zmiana
    współrzędnych wierzchołka jest natychmiast widoczna we wszystkich
    krawędziach, bez konieczności synchronizacji.
  \item \textbf{Czyszczenie grafu} --- metoda \texttt{Graph.clear()}
    usuwa wszystkie referencje z mapy i listy, co umożliwia GC
    zwolnienie pamięci po starych obiektach.
  \item \textbf{Pliki tymczasowe} --- pliki tworzone przez aplikację
    (np. przy czyszczeniu komentarzy) są oznaczane przez
    \texttt{deleteOnExit()}, co zapewnia ich usunięcie przy zamykaniu JVM
    nawet w przypadku nieoczekiwanego zakończenia programu.
  \item \textbf{Strumienie} --- \texttt{Scanner} i \texttt{PrintWriter}
    są używane w blokach \texttt{try-with-resources}, co gwarantuje
    automatyczne zamknięcie strumieni po zakończeniu operacji.
  \item \textbf{Procesy zewnętrzne} --- proces \texttt{graf\_planarny}
    jest uruchamiany przez \texttt{ProcessBuilder}; metoda
    \texttt{waitFor()} blokuje wątek do zakończenia procesu,
    zapobiegając wyciekom procesów.
\end{itemize}

% ============================================================
\section{Kompilacja i uruchamianie}
% ============================================================

\subsection{Wymagania}

\begin{itemize}
  \item Java SE 8 lub nowsza,
  \item skompilowany program \texttt{graf\_planarny} (C) umieszczony
    w podkatalogu \texttt{core/} katalogu roboczego aplikacji,
  \item system operacyjny z graficznym środowiskiem okienkowym,
  \item (opcjonalnie) narzędzie \texttt{make} do uproszczonej kompilacji.
\end{itemize}

\subsection{Metoda uproszczona (zalecana)}

W głównym katalogu projektu znajduje się plik \texttt{Makefile},
który automatyzuje proces kompilacji i uruchomienia. Aby zbudować
i wystartować aplikację, wystarczy wykonać:

\begin{cmdbox}
make
\end{cmdbox}

Polecenie to kompiluje wszystkie pliki \texttt{.java} i natychmiast
uruchamia GUI. Jeśli silnik C w folderze \texttt{core/} nie był jeszcze
budowany, należy uprzednio wykonać:

\begin{cmdbox}
cd core\\
make\\
cd ..
\end{cmdbox}

\subsection{Kompilacja ręczna}

W przypadku braku narzędzia \texttt{make} można skompilować projekt
bezpośrednio:

\begin{cmdbox}
\# Kompilacja wszystkich plików źródłowych:\\
javac -encoding UTF-8 gui/*.java
\end{cmdbox}

W wyniku powstają pliki \texttt{.class} w katalogu \texttt{gui/}.
Aplikację uruchamia się poleceniem:

\begin{cmdbox}
java -cp gui Main
\end{cmdbox}

\subsection{Struktura katalogów wdrożeniowych}

\begin{filebox}
katalog\_projektu/\\
\quad core/\\
\quad \quad graf\_planarny \quad\# (Linux/macOS)\\
\quad \quad graf\_planarny.exe \quad\# (Windows)\\
\quad gui/\\
\quad \quad *.java \quad\# pliki źródłowe\\
\quad \quad *.class \quad\# skompilowane klasy Java\\
\quad Makefile
\end{filebox}

Program C musi mieć prawo do wykonania (na systemach Unix:
\texttt{chmod +x core/graf\_planarny}). Aplikacja Java nie przeszukuje
innych ścieżek --- oczekuje pliku wykonywalnego dokładnie w podkatalogu
\texttt{core/} bieżącego katalogu roboczego.

\end{document}